% META: {"id":"TCOMPL-AIR-EX-002","chapitre":"TCOMPL-CALCULS-AIRES","type_objet":"exercice","capacites":["TCOMPL-AIR-C3"],"capacites_codes":["C3"],"methodes":[],"parcours":1,"competences":["calculer"],"duree_min":15,"mode_creation":"ex_nihilo","sources_inspiration":[],"fichier_tex":"chapitres/TCOMPL-CALCULS-AIRES/exercices/TCOMPL-AIR-EX-002.tex","corrige_tex":"chapitres/TCOMPL-CALCULS-AIRES/corriges/TCOMPL-AIR-CO-002.tex","status":"approved","origin":"generated","status_history":["generated","audited","accepted","approved"],"release_acceptance":"RELEASE_OWNER_BATCH_ACCEPTANCE","acceptance_closure_digest":"sha256:9b3ccf9a81c5520fbb7e03b7b2d3e7bf2057a02834b6d908bf3be5f49f3b553f","acceptance_receipt_digest":"sha256:4fad86f0282e0fa4e0419cca1ad6379ae7af5a280c04a4a90880f90a1c044242"}
% BEGIN-VERIFY
% from sympy import symbols, integrate, Rational
% x = symbols('x')
% f = 3*x**2 - 2*x + 1
% I = integrate(f, (x, 0, 2))
% assert I == 6
% moyenne = I / (2 - 0)
% assert moyenne == 3
% END-VERIFY
\begin{exercice}{TCOMPL-AIR-EX-002}{1}{15}
Soit $f(x) = 3x^2-2x+1$.
\begin{enumerate}
  \item Calculer $\displaystyle\int_0^2 f(x)\,\mathrm{d}x$.
  \item En déduire la valeur moyenne de $f$ sur $[0,2]$.
\end{enumerate}
\end{exercice}
